\newpage
\section{Computational and Python Coding Methods}
\label{sec:computational-methods}

\subsection{Model and assumptions}

Our program calculates what happens to the satellite right after one hit:
the impulse it gets, its velocity change $\Delta\vect V$, its spin
$\omega$, the fragment's velocity after the hit, and where the kinetic
energy goes. We do not step through time. Instead, the program plugs the
inputs into the exact formulas from Section~\ref{sec:basic-physics}. We
made these assumptions:
\begin{itemize}
  \item \textit{The collision is instant.} Outside forces act for too
    short a time to matter, so $\vect P$ and $L_z$ are conserved
    (Eq.~\eqref{eq:conservation}). We do not model how the contact force
    changes during the hit.
  \item \textit{Two dimensions.} Everything moves in the orbital plane and
    the satellite only spins about $z$.
  \item \textit{Rigid satellite, point-mass fragment.} The satellite does
    not bend or break. The fragment's size is only used to report the
    diameter of a sphere with the same mass, $d=(6m/\pi\rho)^{1/3}$.
  \item \textit{Only two ideal types of collision.} Elastic with no
    friction (Eqs.~\eqref{eq:impulse}--\eqref{eq:J-elastic}), or completely
    inelastic (Eqs.~\eqref{eq:rc}--\eqref{eq:omega-inelastic}). A real
    fragment at these speeds melts or vaporises, so the real answer is
    somewhere between the two.
  \item \textit{The orbit is only used for scale.} We only report
    $|\Delta\vect V|/v_{\mathrm{orb}}$; we do not calculate the new orbit.
\end{itemize}

\subsection{Algorithm}

For each hit, the program does these steps.
\begin{enumerate}
  \item Make the unit vectors $\hat{\vect n}=(\cos\alpha,\sin\alpha)$ from
    the panel angle $\alpha$ and $\hat{\vect t}=\hat{\vect z}\times\hat{\vect n}$,
    then find $\vect v_{d0}$ from Eq.~\eqref{eq:vd0}.
  \item \textit{Elastic:} find $b=r_xn_y-r_yn_x$, $k$ from Eq.~\eqref{eq:k}
    and $J$ from Eq.~\eqref{eq:J-elastic}. Then
    $\Delta\vect V=-J\hat{\vect n}/M$, $\omega=-Jb/I$ and
    $\vect v_{d1}=\vect v_{d0}+J\hat{\vect n}/m$ (Eq.~\eqref{eq:impulse}).
  \item \textit{Completely inelastic:} find $\vect r_c$, $\vect v_{CM}$,
    $I_c$ and $\omega$ from Eqs.~\eqref{eq:rc}--\eqref{eq:omega-inelastic}.
    The old CoM of the satellite then moves with
    $\Delta\vect V=\vect v_{CM}+\omega\hat{\vect z}\times(-\vect r_c)$, and
    the stuck fragment moves with
    $\vect v_{d1}=\vect v_{CM}+\omega\hat{\vect z}\times(\vect r-\vect r_c)$.
  \item Find $K_0$, $K_{d1}$, the satellite's kinetic energy of moving and
    turning, and $E_{\mathrm{diss}}$ (Eq.~\eqref{eq:ediss}). The impulse
    on the satellite is reported as $M|\Delta\vect V|$.
  \item Find the speed of a circular orbit. Gravity is the centripetal
    force, $GM_EM/(R_E+h)^2=Mv_{\mathrm{orb}}^2/(R_E+h)$, so
    \begin{equation}
      v_{\mathrm{orb}}=\sqrt{\mu_E/(R_E+h)},
      \label{eq:vorb}
    \end{equation}
    where $\mu_E=GM_E$, $R_E$ is Earth's radius and $h$ is the altitude.
    Then report $|\Delta\vect V|/v_{\mathrm{orb}}$.
  \item Check the conservation laws. The program calculates $\vect P$ and
    $L_z$ (Eq.~\eqref{eq:conservation}) before and after, and the energy
    balance, and reports the relative errors
    \begin{equation}
      \varepsilon_P=\frac{|\vect P_1-\vect P_0|}{|\vect P_0|},\quad
      \varepsilon_L=\frac{|L_{z1}-L_{z0}|}{|L_{z0}|},\quad
      \varepsilon_E=\frac{|K_0-K_{d1}-K_T-K_R-E_{\mathrm{diss}}|}{K_0}.
      \label{eq:residuals}
    \end{equation}
    For a hit through the CoM, $L_{z0}=0$, so $\varepsilon_L$ is divided
    by $|\vect P_0|\,|\vect r|$ instead.
\end{enumerate}
The energy check is a real test. $E_{\mathrm{diss}}$ is calculated on its
own in the inelastic case and is set to zero in the elastic case, so
$\varepsilon_E$ is only zero if the impulses in
Eqs.~\eqref{eq:impulse}--\eqref{eq:omega-inelastic} are correct. For the
mass-ratio sweep (head-on, $\vect r=\vect 0$), the program also compares
its answers with the 1-D formulas in Section~\ref{sec:basic-physics} and
with $v=m_1v_1/(m_1+m_2)$ for the completely inelastic case.

\subsection{Python implementation}

The code is written in Python 3.12 and is run with
\texttt{uv run main.py}. It has five files.
\texttt{input\_parser.py} reads the inputs, changes angles to radians and
stops the program if a value makes no physical sense (mass, moment of
inertia, speed or altitude not positive, unknown collision type, or impact
angle outside $[0^\circ,90^\circ)$). \texttt{physics\_model.py} has the
formulas (Eqs.~\eqref{eq:vd0}--\eqref{eq:ediss} and Eq.~\eqref{eq:vorb})
and the conservation check. \texttt{solvers.py} runs the scenarios, the
parameter sweeps and the limiting cases, \texttt{plotting.py} draws the
figures, and \texttt{main.py} runs everything and saves the tables. The physics functions work on whole NumPy arrays,
so a full sweep is done in one call instead of a loop. We use pandas for
the output tables and Matplotlib for the figures. No physical number is
typed into the code; all of them come from the input files.

\subsection{Input and output design}

Each scenario is one JSON file with four parts: \textit{satellite} ($M$,
$I$, impact point $\vect r$, panel angle $\alpha$), \textit{debris} ($m$,
material, $v_0$, $\theta$), \textit{contact} (\texttt{elastic} or
\texttt{completely\_inelastic}) and \textit{orbit} ($h$). Every key name
includes its unit (for example \texttt{mass\_kg}). A CSV file lists the
density of each material, and a JSON file stores $\mu_E$ and $R_E$.
Another JSON file sets up the parameter study: a $71\times57$ grid of
debris masses ($10^{-6}$--$10\unit{kg}$, log spacing) and speeds
($1$--$15\unit{km\,s^{-1}}$), 179 impact angles from $0^\circ$ to
$89^\circ$ for each collision type, 241 mass ratios $m/M$ from $10^{-6}$
to $100$ for each type, and eight limiting cases.

For each scenario, the program saves \texttt{summary.csv} (all results
and errors) and \texttt{conservation.csv} (the values of $\vect P$, $L_z$
and energy before and after). It also saves combined tables for all
scenarios and one CSV file per sweep. We designed each figure to show one
idea: bar charts of momentum (normal and tangential parts) and energy
before and after the hit, a colour map of $|\Delta\vect V|$ against mass
and speed, line graphs of $|\Delta\vect V|$ and $|\omega|$ against angle
and mass ratio, and a bar chart of the conservation errors. We use log
scales when a quantity covers many orders of magnitude.
